\section{Discussion}

The central technical idea is to evaluate legal reasoning at the cluster level
rather than scoring every response independently. If Dasha can reliably group
answers by reasoning path, a judge can score representative centroids and then
project those scores to all cluster members. This creates an interpretable
middle layer between raw model generations and aggregate model rankings.

The reported evidence supports the feasibility of that design on a small
live natural-response batch. The live run shows that the source-to-ranking path
works with real model outputs and that Dasha can separate more than one observed
legal-reasoning path without preassigned response labels. The controlled
500-response regression test shows that, once reasoning signatures are
available, Dasha can preserve eight legally distinct Statute-of-Frauds reasoning
paths at the intended compression scale. These two evidence layers answer
different questions: the live run tests discovery from natural answers, while
the regression suite tests scale, metrics, and bookkeeping under known
archetypes.

The stage-level assessment is equally important. Frank and Karthic are not
validated merely by emitting JSON. They are validated by whether the generated
question, gold answer, variations, and rubric rows preserve the legally
operative facts and match the methodology described in the instruction context.
Likewise, Dasha is not validated merely by reporting a cluster count; its
clusters must be interpretable as different legal theories. Judge is not
validated merely by producing a scalar score; it must apply the Karthic rows to
the centroid and leave row-level rationales that can be reviewed. This is why
the run bundle and paper appendix include concrete artifacts rather than only
aggregate metrics.

The broader transfer claim is deliberately modest. The architecture is designed
so that the case source and instruction stream determine the doctrine-specific
content, while the pipeline mechanics remain stable: source extraction,
question generation, perturbation-track generation, rubric construction,
natural response generation, reasoning clustering, centroid judging, and
escalation. That gives a plausible path to other areas of law, but it is not
yet broad empirical validation. The next stage is to run the same stage-level
and perturbation review protocol on held-out cases and additional doctrine
families without changing the pipeline in between.

The remaining risk is not only the cluster container itself. Dasha depends on an
LLM semantic-extraction step; if that step assigns incorrect normalized
signatures, deterministic grouping will preserve the wrong abstraction. The
judge is also a single rubric-conditioned model in the reported run, so the
reported model rankings should be read as pipeline outputs rather than stable
comparative claims about the model families. The next methodological improvement is a
larger held-out natural-response suite with repeated signature extraction, judge
ensembling, rubric ablations, and adversarial signature audits.
